% ============================================================
%  Guía de aplicación: del dato crudo al hallazgo
%  Semana 6 · Clase 12: MasterClass (Recapitulación del Módulo 2)
%  Curso: Análisis y Visualización de Datos para la Toma de Decisiones
%  Compilar con: xelatex guia-aplicacion-modulo2.tex (dos veces)
% ============================================================
\documentclass[11pt]{article}

\usepackage[a4paper,margin=2cm]{geometry}
\usepackage{fontspec}
\usepackage{amssymb}
\usepackage[spanish,es-noshorthands,es-tabla]{babel}
\usepackage{xcolor}
\usepackage{graphicx}
\usepackage{booktabs}
\usepackage{array}
\usepackage{tabularx}
\usepackage{enumitem}
\usepackage{fancyvrb}
\usepackage{needspace}
\usepackage{tcolorbox}
\usepackage{tikz}
\usepackage{fancyhdr}
\usepackage{hyperref}

\tcbuselibrary{skins,breakable}
\usetikzlibrary{arrows.meta,positioning,calc,shapes.geometric,babel}

% ---------------- Tipografía (misma lógica que la plantilla de slides) ----------------
\IfFontExistsTF{Aptos}{\setmainfont{Aptos}}{\setmainfont{Lato}}
\setmonofont{DejaVu Sans Mono}[Scale=0.82]

% ---------------- Paleta Esumer ----------------
\definecolor{esumerNavy}{HTML}{262E66}
\definecolor{esumerNavyDark}{HTML}{1E4093}
\definecolor{esumerTeal}{HTML}{0CB3B6}
\definecolor{esumerSky}{HTML}{26A1DB}
\definecolor{esumerOrange}{HTML}{EF7C2C}
\definecolor{esumerGreen}{HTML}{94C122}
\definecolor{esumerMagenta}{HTML}{AF348B}
\definecolor{esumerGray}{HTML}{E8E8E8}
\definecolor{esumerBg}{HTML}{F7F8FA}

\hypersetup{colorlinks=true, linkcolor=esumerNavyDark, urlcolor=esumerNavyDark}
\graphicspath{{../../../plantilla/assets/}}

% ---------------- Encabezado / pie de página ----------------
\pagestyle{fancy}
\fancyhf{}
\renewcommand{\headrulewidth}{0pt}
\fancyfoot[L]{\small\color{esumerNavyDark} Guía de aplicación del Módulo 2 \textperiodcentered\ Clase 12}
\fancyfoot[R]{\small\color{esumerNavyDark} Cristian A. Jiménez C. \quad\textperiodcentered\quad \thepage}
\setlength{\parindent}{0pt}
\setlength{\parskip}{4pt}

% ---------------- Código ----------------
\fvset{fontsize=\small}
\newtcolorbox{codigo}[1][]{
  colback=esumerBg, colframe=esumerNavy!35, boxrule=0.5pt, arc=1mm,
  left=6pt, right=6pt, top=2pt, bottom=2pt, fonttitle=\bfseries\small,
  coltitle=esumerNavy, colbacktitle=esumerGray, #1
}

% ---------------- Cajas reutilizables ----------------
\newtcolorbox{recuerda}[1][Lo que pasó en el curso]{
  colback=esumerBg, colframe=esumerTeal, boxrule=0.6pt, arc=1mm,
  fonttitle=\bfseries\small, coltitle=white, colbacktitle=esumerTeal,
  title=#1, left=6pt, right=6pt, top=3pt, bottom=3pt
}
\newtcolorbox{cuidado}[1][Cuidado]{
  colback=white, colframe=esumerOrange, boxrule=0.6pt, arc=1mm,
  fonttitle=\bfseries\small, coltitle=white, colbacktitle=esumerOrange,
  title=#1, left=6pt, right=6pt, top=3pt, bottom=3pt
}

% ---------------- Encabezado de cada paso ----------------
\newcommand{\paso}[3]{%
  \clearpage\noindent
  \begin{tikzpicture}[baseline=(n.base)]
    \node[circle, fill=esumerOrange, text=white, font=\bfseries\large,
          minimum size=0.85cm, inner sep=0pt] (n) {#1};
  \end{tikzpicture}\hspace{8pt}%
  {\Large\bfseries\color{esumerNavy} #2}\par
  \vspace{-2pt}{\color{esumerOrange}\rule{\textwidth}{1pt}}\par
  \vspace{2pt}
  {\color{esumerNavyDark}\textbf{Pregunta guía:}} \textit{#3}\par
  \vspace{2pt}
}

% ---------------- Tabla de funciones de cada paso ----------------
\newenvironment{funciones}{%
  \par\small\noindent
  \tabularx{\textwidth}{@{}>{\raggedright\arraybackslash}p{5.0cm}
    >{\raggedright\arraybackslash}X >{\centering\arraybackslash}p{1.2cm}@{}}
  \toprule
  \textbf{Para\ldots} & \textbf{Función o método} & \textbf{Clase}\\
  \midrule}{%
  \bottomrule\endtabularx\par}

% ---------------- Lista de chequeo ----------------
\newlist{checklist}{itemize}{1}
\setlist[checklist]{label=$\square$, itemsep=1pt, topsep=2pt, leftmargin=1.4em}

% ---------------- Estilos TikZ comunes ----------------
\tikzset{
  every node/.append style={execute at begin node={\hyphenpenalty=10000\exhyphenpenalty=10000\relax}},
  flecha/.style={-{Stealth[length=2.4mm]}, line width=0.9pt, color=esumerNavy},
  caja/.style={draw=esumerNavy, line width=0.7pt, rounded corners=3pt,
               fill=white, align=center, inner sep=4pt},
  fn/.style={font=\scriptsize\ttfamily, text=esumerNavyDark, align=center},
  num/.style={circle, fill=esumerOrange, text=white, font=\bfseries\small,
              minimum size=0.6cm, inner sep=0pt},
}

% ============================================================
\begin{document}

% ---------------- Encabezado ----------------
\begin{minipage}[c]{0.18\textwidth}
  \includegraphics[height=1.15cm,width=\linewidth,keepaspectratio]{esumer-logo.png}
\end{minipage}%
\begin{minipage}[c]{0.80\textwidth}
  \raggedleft
  {\Large\bfseries\color{esumerNavy} Guía de aplicación: del dato crudo al hallazgo}\\[3pt]
  {\small\color{esumerNavyDark} Semana 6 \textperiodcentered\ Clase 12 (MasterClass del Módulo 2) \textperiodcentered\ Análisis y Visualización de Datos para la Toma de Decisiones}
\end{minipage}

\vspace{4pt}
{\color{esumerOrange}\rule{\textwidth}{2pt}}
\vspace{6pt}

\textbf{Para qué es esta guía.} En el Módulo 2 aprendiste cada herramienta
por separado, siempre sobre los datos del gimnasio. Hoy las juntas, en orden,
sobre \textbf{tus propios datos}. La guía te lleva paso a paso: qué
pregunta responde cada etapa, un diagrama de lo que pasa con la tabla, las
funciones que necesitas y una lista de chequeo para saber cuándo puedes
pasar a la siguiente.

\textbf{Al terminar la sesión deberías tener:}
\begin{enumerate}[itemsep=1pt, topsep=2pt]
  \item Un notebook de Colab que va de tu archivo crudo a una tabla limpia y
    enriquecida, sin editar nunca el archivo original.
  \item Una tabla resumen (\texttt{groupby}) que responde al menos una de
    tus preguntas de negocio.
  \item Dos o tres gráficos, cada uno con su hallazgo escrito.
  \item Una \textbf{bitácora de decisiones} (plantilla en la última página):
    qué problemas encontraste y qué decidiste hacer con cada uno.
\end{enumerate}

\textbf{Cómo leer los ejemplos de código.} Usamos nombres genéricos que
debes reemplazar por los tuyos:

\begin{tcolorbox}[colback=esumerBg, colframe=esumerGray, boxrule=0.5pt, arc=1mm, left=6pt, right=6pt, top=3pt, bottom=3pt]
\small
\begin{tabular}{@{}>{\raggedright\arraybackslash}p{3.3cm}p{12.6cm}@{}}
  \texttt{df} & tu tabla tal como la cargaste (cruda). \\
  \texttt{limpio} & una copia de \texttt{df} sobre la que haces todos los cambios. \\
  \texttt{"ciudad"}, \texttt{"fecha"}, \texttt{"monto"}\ldots & nombres de columna de ejemplo: cámbialos por los nombres reales de tu dataset. \\
  \texttt{"id\_cliente"} & la columna que identifica cada fila (cliente, venta, producto\ldots), si tu dataset tiene una. \\
\end{tabular}
\end{tcolorbox}

\textbf{Agenda sugerida para las 3 horas}

\begin{tabular}{@{}>{\raggedright\arraybackslash}p{1.4cm}>{\raggedright\arraybackslash}p{7.3cm}>{\raggedright\arraybackslash}p{7.2cm}@{}}
\toprule
\textbf{Hora} & \textbf{Pasos} & \textbf{Producto al final de la hora}\\
\midrule
Hora 1 & 0. Preguntas \textperiodcentered\ 1. Cargar \textperiodcentered\ 2. Diagnosticar \textperiodcentered\ 3. Limpiar & Tabla \texttt{limpio} sin duplicados, con tipos correctos y nulos tratados.\\
Hora 2 & 4. Transformar \textperiodcentered\ 5. Integrar \textperiodcentered\ 6. Agrupar & Columnas nuevas y una tabla resumen por grupo.\\
Hora 3 & 7. Visualizar y comunicar & Dos o tres gráficos con hallazgo, y la bitácora completa.\\
\bottomrule
\end{tabular}

\vspace{6pt}
La referencia completa de funciones del curso está en
\texttt{recursos/referencia-funciones-metodos.pdf}: tenla abierta mientras
trabajas.

% ============================================================
\clearpage
\section*{\color{esumerNavy}El mapa completo}

Estos son los ocho pasos de la guía. La columna \textbf{Clase} de cada tabla
de funciones te dice dónde lo viste por primera vez, por si necesitas volver
al notebook de esa clase.

\vspace{4pt}
\begin{center}
\begin{tikzpicture}[
  etapa/.style={draw=esumerNavy, line width=0.8pt, rounded corners=3pt,
                fill=white, minimum width=3.74cm, minimum height=2.55cm,
                text width=3.5cm, align=center, inner sep=3pt},
]
  % Fila 1
  \node[etapa] (p0) at (2.2,0)
    {{\bfseries\small Definir\\preguntas}\\[3pt]
     {\scriptsize ¿Qué decisión quiero apoyar con estos datos?}};
  \node[etapa] (p1) at (6.5,0)
    {{\bfseries\small Cargar\\y conocer}\\[3pt]
     {\scriptsize\ttfamily read\_csv\\shape \textperiodcentered\ head\\info \textperiodcentered\ describe}};
  \node[etapa] (p2) at (10.8,0)
    {{\bfseries\small Diagnosticar\\la calidad}\\[3pt]
     {\scriptsize\ttfamily isna \textperiodcentered\ duplicated\\dtypes\\value\_counts}};
  \node[etapa] (p3) at (15.1,0)
    {{\bfseries\small Limpiar}\\[3pt]
     {\scriptsize\ttfamily drop\_duplicates\\str.strip \textperiodcentered\ replace\\to\_numeric \textperiodcentered\ fillna}};
  % Fila 2
  \node[etapa] (p4) at (2.2,-3.7)
    {{\bfseries\small Transformar\\y enriquecer}\\[3pt]
     {\scriptsize\ttfamily pd.cut \textperiodcentered\ map\\dt.to\_period\\columnas calculadas}};
  \node[etapa] (p5) at (6.5,-3.7)
    {{\bfseries\small Integrar\\fuentes}\\[1pt]{\scriptsize (si tienes más de una)}\\[2pt]
     {\scriptsize\ttfamily merge \textperiodcentered\ concat\\validate}};
  \node[etapa] (p6) at (10.8,-3.7)
    {{\bfseries\small Agrupar\\y resumir}\\[3pt]
     {\scriptsize\ttfamily groupby \textperiodcentered\ agg\\transform\\crosstab}};
  \node[etapa] (p7) at (15.1,-3.7)
    {{\bfseries\small Visualizar\\y comunicar}\\[3pt]
     {\scriptsize\ttfamily .plot \textperiodcentered\ .plot.bar\\.plot.hist}\\[1pt]
     {\scriptsize título = conclusión}};
  % Números
  \foreach \i in {0,...,7} {\node[num] at (p\i.north west) {\i};}
  % Flechas
  \draw[flecha] (p0) -- (p1);
  \draw[flecha] (p1) -- (p2);
  \draw[flecha] (p2) -- (p3);
  \draw[flecha] (p3.south) -- ++(0,-0.5) -| (p4.north);
  \draw[flecha] (p4) -- (p5);
  \draw[flecha] (p5) -- (p6);
  \draw[flecha] (p6) -- (p7);
  % Banda de validación continua
  \node[draw=esumerTeal, dashed, line width=0.8pt, rounded corners=3pt,
        fill=esumerBg, text width=16.1cm, align=left, inner sep=5pt]
    at (8.65,-6.0)
    {\small\textbf{\color{esumerNavy}Validación continua.} Después de
     \emph{cada} paso cuenta filas (\texttt{len}), revisa nulos
     (\texttt{isna().sum()}) y mira los valores (\texttt{value\_counts()}).
     Si un número no cuadra, no sigas: vuelve al paso anterior y entiende
     por qué.};
\end{tikzpicture}
\end{center}

\vspace{6pt}
Tres reglas que atraviesan todos los pasos (las viste en la Clase 8):

\begin{enumerate}[itemsep=2pt, topsep=2pt]
  \item \textbf{Nunca edites el archivo crudo.} Todo parte de
    \texttt{pd.read\_csv(...)}; si te equivocas, vuelves a leerlo y no
    perdiste nada.
  \item \textbf{Trabaja sobre una copia} (\texttt{limpio = df.copy()}) y,
    cuando corrijas una columna, compárala con la original antes de
    reemplazarla.
  \item \textbf{Nada que cambie cada vez que lo ejecutas.} Si calculas
    antigüedades, fija una fecha de referencia
    (\texttt{pd.Timestamp("2026-09-30")}), no uses la fecha de hoy.
\end{enumerate}

% ============================================================
\paso{0}{Definir tus preguntas}{¿Qué decisión quiero apoyar, y qué tengo que saber para tomarla?}

Antes de abrir el notebook, escribe dos o tres preguntas de negocio. Cada
una te dice qué columnas necesitas y, desde ya, qué tipo de gráfico vas a
usar al final (Clase 11). Una pregunta concreta también te dice qué
limpiar primero: no hace falta dejar perfecta una columna que no vas a usar.

\begin{center}
\begin{tikzpicture}[
  eslabon/.style={caja, minimum width=3.8cm, text width=3.5cm, minimum height=2.3cm},
  tit/.style={font=\bfseries\small, text=white, fill=esumerNavy, rounded corners=2pt,
              minimum width=3.8cm, inner sep=3pt},
]
  \foreach \x/\tit in {2.2/Decisión, 6.5/Pregunta, 10.8/Columnas, 15.1/Tipo de pregunta}
    {\node[tit] at (\x,1.4) {\tit};}
  \node[eslabon] (a) at (2.2,0) {\small ¿En qué ciudad reforzamos la campaña de renovación?};
  \node[eslabon] (b) at (6.5,0) {\small ¿Qué ciudad concentra más socios activos?};
  \node[eslabon] (c) at (10.8,0) {\small\texttt{ciudad}\\\texttt{activo}};
  \node[eslabon] (d) at (15.1,0) {\small\textbf{Comparación}\\pocas categorías\\$\Rightarrow$ barras};
  \draw[flecha] (a) -- (b);
  \draw[flecha] (b) -- (c);
  \draw[flecha] (c) -- (d);
\end{tikzpicture}

{\footnotesize\itshape Ejemplo con los datos del gimnasio. Los cuatro tipos
de pregunta son los de la Clase 11: comparación, composición, distribución y relación.}
\end{center}

\textbf{Escribe aquí tus preguntas:}

\renewcommand{\arraystretch}{2.1}
\begin{tabularx}{\textwidth}{@{}>{\centering\arraybackslash}p{0.5cm} >{\raggedright\arraybackslash}X >{\raggedright\arraybackslash}p{3.6cm} >{\raggedright\arraybackslash}p{3.0cm}@{}}
\toprule
& \textbf{Pregunta de negocio} & \textbf{Columnas que necesito} & \textbf{Tipo de pregunta}\\
\midrule
1 & & & \\ \hline
2 & & & \\ \hline
3 & & & \\
\bottomrule
\end{tabularx}\par
\renewcommand{\arraystretch}{1.15}

\begin{checklist}
  \item Cada pregunta se puede responder con columnas que existen en mis datos.
  \item Cada pregunta tiene asignado uno de los cuatro tipos.
\end{checklist}

% ============================================================
\paso{1}{Cargar y conocer tus datos}{¿Qué tengo en las manos: cuántas filas, qué columnas y de qué tipo?}

La función de lectura depende del formato del archivo. Todas devuelven lo
mismo: un \texttt{DataFrame}. A partir de ahí, cuatro órdenes te dan la
primera foto de la tabla.

\begin{center}
\begin{tikzpicture}[
  ext/.style={caja, minimum width=1.9cm, font=\small\bfseries},
  lect/.style={caja, minimum width=4.9cm, text width=4.7cm, font=\scriptsize\ttfamily},
]
  \node[caja, minimum width=2.3cm, minimum height=1.6cm, fill=esumerBg] (arch) at (1.2,0)
    {\small Tu\\archivo};
  \node[ext] (e1) at (3.9,1.1) {.csv};
  \node[ext] (e2) at (3.9,0) {.xlsx};
  \node[ext] (e3) at (3.9,-1.1) {.db};
  \node[lect] (l1) at (8.0,1.1) {pd.read\_csv("archivo.csv")};
  \node[lect] (l2) at (8.0,0) {pd.read\_excel("archivo.xlsx")};
  \node[lect] (l3) at (8.0,-1.1) {pd.read\_sql\_query(consulta, con)};
  \node[circle, draw=esumerNavy, fill=esumerTeal, text=white, line width=0.8pt,
        minimum size=1.1cm, font=\bfseries\ttfamily] (df) at (11.55,0) {df};
  \node[caja, text width=4.0cm, fill=esumerBg] (vis) at (14.85,0)
    {\small\bfseries Primera foto\\[2pt]
     \scriptsize\ttfamily df.shape\\df.head() \textperiodcentered\ df.sample(5)\\df.info()\\df.describe(include="all")};
  \foreach \i in {1,2,3} {
    \draw[flecha] (arch.east) -- (e\i.west);
    \draw[flecha] (e\i) -- (l\i);
    \draw[flecha] (l\i.east) -- (df);
  }
  \draw[flecha] (df) -- (vis);
\end{tikzpicture}
\end{center}

\begin{funciones}
Leer un CSV o un Excel & \texttt{pd.read\_csv("archivo.csv")}, \texttt{pd.read\_excel("archivo.xlsx")} & 4, 7\\
Leer una tabla de una base SQLite & \texttt{pd.read\_sql\_query("SELECT * FROM tabla", con)} & 5\\
Saber el tamaño (filas, columnas) & \texttt{df.shape} & 7\\
Ver filas de muestra & \texttt{df.head()}, \texttt{df.sample(5)} & 7\\
Tipos y valores no nulos por columna & \texttt{df.info()} & 7\\
Resumen estadístico (también del texto) & \texttt{df.describe(include="all")} & 7\\
\end{funciones}

\begin{codigo}[title=En tu notebook]
\begin{Verbatim}
import pandas as pd

df = pd.read_csv("mis_datos.csv")
df.shape
df.head()
df.info()
\end{Verbatim}
\end{codigo}

\begin{cuidado}[Si el archivo no se lee bien]
\small Archivos exportados desde Excel en español suelen usar punto y coma
como separador y coma decimal; si todas las columnas aparecen pegadas en
una sola, prueba \texttt{pd.read\_csv("archivo.csv", sep=";", decimal=",")}.
Si las tildes salen como símbolos raros, agrega \texttt{encoding="latin-1"}.
\end{cuidado}

\begin{checklist}
  \item Sé cuántas filas y columnas tiene mi tabla, y qué representa cada fila.
  \item Revisé \texttt{df.info()}: anoté las columnas cuyo tipo no es el que esperaba.
\end{checklist}

% ============================================================
\paso{2}{Diagnosticar la calidad}{¿Qué problemas tiene la tabla antes de confiar en ella?}

Busca, columna por columna, los cuatro problemas de la Clase 7. El diagrama
muestra cómo se ven en una tabla y qué función los delata.

\begin{center}
\begin{tikzpicture}[
  celda/.style={draw=esumerGray!70!black, line width=0.4pt, minimum height=0.55cm,
                inner sep=1pt, font=\scriptsize},
  cab/.style={celda, fill=esumerNavy, text=white, font=\scriptsize\bfseries},
]
  % columnas: centro / ancho
  \def\cA{0.6}  \def\wA{1.2}
  \def\cB{2.35} \def\wB{2.3}
  \def\cC{4.4}  \def\wC{1.8}
  \def\cD{6.6}  \def\wD{2.6}
  % encabezado
  \node[cab, minimum width=\wA cm] at (\cA,0) {id};
  \node[cab, minimum width=\wB cm] at (\cB,0) {ciudad};
  \node[cab, minimum width=\wC cm] at (\cC,0) {edad};
  \node[cab, minimum width=\wD cm] at (\cD,0) {telefono};
  % fila 1
  \node[celda, minimum width=\wA cm] at (\cA,-0.55) {1};
  \node[celda, minimum width=\wB cm] at (\cB,-0.55) {Medellín};
  \node[celda, minimum width=\wC cm] at (\cC,-0.55) {34};
  \node[celda, minimum width=\wD cm] at (\cD,-0.55) {3104451122};
  % fila 2
  \node[celda, minimum width=\wA cm] at (\cA,-1.1) {2};
  \node[celda, minimum width=\wB cm, fill=esumerGreen!35] at (\cB,-1.1) {MEDELLIN};
  \node[celda, minimum width=\wC cm, fill=esumerSky!30] at (\cC,-1.1) {treinta};
  \node[celda, minimum width=\wD cm, fill=esumerSky!30] at (\cD,-1.1) {3.031046e+09};
  % fila 3
  \node[celda, minimum width=\wA cm] at (\cA,-1.65) {3};
  \node[celda, minimum width=\wB cm] at (\cB,-1.65) {Bogotá};
  \node[celda, minimum width=\wC cm, fill=esumerOrange!35] at (\cC,-1.65) {NaN};
  \node[celda, minimum width=\wD cm] at (\cD,-1.65) {3012267788};
  % fila 4
  \node[celda, minimum width=\wA cm] at (\cA,-2.2) {4};
  \node[celda, minimum width=\wB cm, fill=esumerGreen!35] at (\cB,-2.2) {\texttt{" Medellín"}};
  \node[celda, minimum width=\wC cm] at (\cC,-2.2) {41};
  \node[celda, minimum width=\wD cm, fill=esumerOrange!35] at (\cD,-2.2) {NaN};
  % fila 5 (duplicado exacto de la 4)
  \node[celda, minimum width=\wA cm, fill=esumerMagenta!25] at (\cA,-2.75) {4};
  \node[celda, minimum width=\wB cm, fill=esumerMagenta!25] at (\cB,-2.75) {\texttt{" Medellín"}};
  \node[celda, minimum width=\wC cm, fill=esumerMagenta!25] at (\cC,-2.75) {41};
  \node[celda, minimum width=\wD cm, fill=esumerMagenta!25] at (\cD,-2.75) {NaN};
  \draw[esumerMagenta, line width=1.2pt, rounded corners=1pt] (0,-3.025) rectangle (7.9,-2.475);
  % leyenda
  \foreach \y/\col/\tit/\txt/\fun in {
      0.05/{esumerOrange!35}/{Nulo}/{una celda vacía (NaN)}/{df.isna().sum()},
      -0.95/{esumerMagenta!25}/{Duplicado}/{una fila idéntica a otra}/{df.duplicated().sum()},
      -1.95/{esumerSky!30}/{Tipo incorrecto}/{texto donde van números}/{df.dtypes \textperiodcentered\ pd.to\_numeric(..., errors="coerce")},
      -2.95/{esumerGreen!35}/{Inconsistencia}/{un valor, varias escrituras}/{df["ciudad"].value\_counts()}} {
    \node[draw=esumerGray!70!black, fill=\col, minimum size=0.42cm, inner sep=0pt] at (8.75,\y) {};
    \node[anchor=north west, text width=7.6cm, inner sep=0pt, align=left] at (9.15,\y+0.2)
      {\small\textbf{\tit}: \txt.\\[-1pt]{\scriptsize\ttfamily\color{esumerNavyDark}\fun}};
  }
\end{tikzpicture}

{\footnotesize\itshape Tabla ilustrativa, con los mismos tipos de problema que encontramos en
\texttt{membresias\_gimnasio.csv} en la Clase 7.}
\end{center}

\begin{funciones}
Contar nulos por columna & \texttt{df.isna().sum()} & 7\\
Ver el porcentaje de nulos por columna & \texttt{(df.isna().mean() * 100).round(1)} & 7\\
Contar filas exactamente repetidas & \texttt{df.duplicated().sum()} & 7\\
Ver si la llave se repite & \texttt{df["id\_cliente"].duplicated().sum()} o \texttt{.is\_unique} & 7, 10\\
Ver el tipo de cada columna & \texttt{df.dtypes} & 7\\
Encontrar texto disfrazado de número & \texttt{pd.to\_numeric(df["edad"], errors="coerce")} & 7\\
Ver las variantes de una columna de texto & \texttt{df["ciudad"].value\_counts()}, \texttt{.unique()} & 7\\
\end{funciones}

\begin{codigo}[title=En tu notebook: el diagnóstico completo]
\begin{Verbatim}
df.isna().sum()                      # nulos por columna
df.duplicated().sum()                # filas repetidas
df.dtypes                            # ¿cada columna tiene el tipo que esperas?

# texto disfrazado de número: se vuelve NaN al convertir, pero no estaba vacío
numero = pd.to_numeric(df["edad"], errors="coerce")
df.loc[numero.isna() & df["edad"].notna(), "edad"]

for columna in ["ciudad", "categoria"]:   # repite con tus columnas de texto
    print(df[columna].value_counts(), "\n")
\end{Verbatim}
\end{codigo}

\begin{recuerda}
\small El archivo crudo del gimnasio tenía 54 filas, de las cuales 4 eran
duplicados exactos; \texttt{ciudad} tenía 12 formas de escribir 4 ciudades
(\texttt{Medellín}, \texttt{Medellin}, \texttt{MEDELLIN},
\texttt{" Medellín"}, \texttt{medellín}\ldots); y \texttt{edad} tenía el
texto \texttt{"treinta"}, que \texttt{isna()} nunca habría mostrado
porque técnicamente no es un nulo.
\end{recuerda}

\begin{checklist}
  \item Para cada columna que uso en mis preguntas, sé si tiene nulos, tipo
    incorrecto o inconsistencias.
  \item Sé si hay filas duplicadas y si mi llave (si la hay) es única.
  \item Anoté cada problema en la bitácora, antes de corregir nada.
\end{checklist}

% ============================================================
\paso{3}{Limpiar}{¿Qué corrijo, en qué orden, y qué hago con lo que no se puede corregir?}

El orden importa. Los duplicados van primero, porque si no, cada conteo y
cada promedio posterior queda inflado por las copias. Los nulos van al
final, porque una conversión de tipo puede crear nulos nuevos
(\texttt{errors="coerce"}).

\begin{center}
\begin{tikzpicture}[
  chip/.style={caja, minimum width=2.45cm, text width=2.15cm, minimum height=1.05cm,
               fill=esumerBg, font=\small\bfseries},
]
  \foreach \i/\x/\tit/\fun in {
      1/1.5/Copia/limpio = df.copy(),
      2/4.3/Duplicados/drop\_duplicates(),
      3/7.1/Texto/str.strip()\\str.lower(),
      4/9.9/Tipos/pd.to\_numeric()\\astype(),
      5/12.7/Nulos/fillna()\\dropna(),
      6/15.5/Verifica/len() \textperiodcentered\ isna()\\value\_counts()} {
    \node[chip] (c\i) at (\x,0) {\tit};
    \node[num, minimum size=0.5cm, font=\bfseries\scriptsize] at (c\i.north west) {\i};
    \node[fn, text width=2.5cm] at (\x,-1.05) {\fun};
  }
  \foreach \i [evaluate=\i as \j using int(\i+1)] in {1,...,5}
    {\draw[flecha] (c\i) -- (c\j);}
\end{tikzpicture}
\end{center}

\textbf{¿Qué hago con los nulos?} No hay una única respuesta correcta: hay una
decisión que tomar, y que debe quedar escrita. Este árbol resume las
decisiones que discutimos en la Clase 8.

\begin{center}
\begin{tikzpicture}[
  preg/.style={caja, fill=esumerNavy, text=white, draw=esumerNavy, font=\small\bfseries},
  hoja/.style={caja, fill=esumerBg, font=\small},
  etq/.style={font=\scriptsize\itshape, text=esumerNavyDark, above=1pt, inner sep=1pt},
]
  % Nivel 0
  \node[preg, text width=6.2cm, anchor=north] (q1) at (8.5,0) {¿La columna es una llave o un identificador?};
  % Nivel 1
  \node[hoja, text width=5.4cm, anchor=north] (h1) at (3.0,-2.2)
    {No inventes valores: revisa el origen. Si no se puede recuperar, elimina la fila.\\
     {\scriptsize\ttfamily dropna(subset=["id\_cliente"])}};
  \node[preg, text width=5.0cm, anchor=north] (q2) at (12.3,-2.2) {¿Qué proporción de la columna está vacía?};
  % Nivel 2
  \node[hoja, text width=5.6cm, anchor=north] (h2) at (6.9,-4.5)
    {Pregúntate si la columna sirve para tus preguntas. Si no la necesitas, no la uses; si la necesitas, documenta la limitación.};
  \node[preg, text width=4.0cm, anchor=north] (q3) at (13.4,-4.5) {¿Qué tipo de dato es?};
  % Nivel 3
  \node[hoja, text width=3.5cm, anchor=north] (h3) at (11.3,-6.4)
    {Imputa con la mediana.\\
     {\scriptsize\ttfamily fillna(s.median())}};
  \node[hoja, text width=3.5cm, anchor=north] (h4) at (15.15,-6.4)
    {Usa una categoría explícita.\\
     {\scriptsize\ttfamily fillna("Sin dato")}};
  % Conexiones: la etiqueta va sobre el tramo horizontal de cada rama
  \draw[flecha] (q1.south) -- (8.5,-1.6) -| (h1.north) node[etq, pos=0.25] {sí};
  \draw[flecha] (q1.south) -- (8.5,-1.6) -| (q2.north) node[etq, pos=0.25] {no};
  \draw[flecha] (q2.south) -- (12.3,-3.9) -| (h2.north) node[etq, pos=0.25] {mucha (cerca de la mitad o más)};
  \draw[flecha] (q2.south) -- (12.3,-3.9) -| (q3.north) node[etq, pos=0.35] {poca};
  \draw[flecha] (q3.south) -- (13.4,-5.8) -| (h3.north) node[etq, pos=0.25] {numérica};
  \draw[flecha] (q3.south) -- (13.4,-5.8) -| (h4.north) node[etq, pos=0.25] {texto};
\end{tikzpicture}
\end{center}

\begin{funciones}
Quitar filas repetidas & \texttt{limpio.drop\_duplicates()} & 8\\
Quitar espacios y unificar mayúsculas & \texttt{.str.strip()}, \texttt{.str.lower()}, \texttt{.str.title()} & 8\\
Unificar variantes que quedan (tildes) & \texttt{.replace(\{"medellin": "medellín"\})} & 8, 10\\
Convertir texto a número & \texttt{pd.to\_numeric(s, errors="coerce")} & 7, 8\\
Convertir a texto (códigos, teléfonos) & \texttt{s.astype(str)} & 8\\
Rellenar nulos & \texttt{s.fillna(s.median())}, \texttt{s.fillna("Sin dato")} & 8\\
Eliminar filas con nulos en columnas clave & \texttt{limpio.dropna(subset=["id\_cliente"])} & 12\\
\end{funciones}

\begin{codigo}[title=En tu notebook]
\begin{Verbatim}
limpio = df.copy()
limpio = limpio.drop_duplicates()

limpio["ciudad_limpia"] = (limpio["ciudad"].str.strip()
                                           .str.lower()
                                           .replace({"medellin": "medellín", "bogota": "bogotá"})
                                           .str.title())
limpio["ciudad_limpia"].value_counts()      # compara con la original antes de reemplazar

limpio["edad"] = pd.to_numeric(limpio["edad"], errors="coerce")
print("edades a imputar:", limpio["edad"].isna().sum())   # anótalo en la bitácora
limpio["edad"] = limpio["edad"].fillna(limpio["edad"].median())

print(len(df), "->", len(limpio))
\end{Verbatim}
\end{codigo}

\begin{checklist}
  \item No quedan duplicados, y sé cuántas filas eliminé.
  \item Cada columna de texto que uso tiene solo valores válidos (revisado con
    \texttt{value\_counts()}).
  \item Las columnas numéricas son números y las fechas son fechas
    (\texttt{limpio.dtypes}).
  \item Para cada columna con nulos tomé una decisión, y escribí cuántos
    valores imputé o eliminé.
\end{checklist}

% ============================================================
\paso{4}{Transformar y enriquecer}{¿Qué columnas nuevas necesito para responder mis preguntas?}

Las columnas originales casi nunca son las que necesitas. Cuatro
transformaciones cubren la mayoría de los casos:

\begin{center}
\begin{tikzpicture}[
  panel/.style={draw=esumerNavy!40, line width=0.6pt, rounded corners=3pt, fill=esumerBg},
  tpanel/.style={anchor=north west, fill=esumerNavy, text=white, font=\scriptsize\bfseries,
                 rounded corners=2pt, inner sep=3pt},
  dato/.style={caja, font=\scriptsize, minimum height=0.85cm},
]
  % ---------- Panel A: columna calculada ----------
  \draw[panel] (0.05,0) rectangle (8.2,3.2);
  \node[tpanel] at (0.15,3.05) {A. Columna calculada};
  \node[dato, minimum width=1.6cm] (a1) at (1.2,1.6) {\textbf{cantidad}\\3};
  \node at (2.35,1.6) {\large$\times$};
  \node[dato, minimum width=1.8cm] (a2) at (3.6,1.6) {\textbf{precio}\\20.000};
  \node at (4.85,1.6) {\large$=$};
  \node[dato, minimum width=1.8cm, fill=esumerTeal!25] (a3) at (6.2,1.6) {\textbf{total}\\60.000};
  \node[fn] at (4.1,0.45) {df["total"] = df["cantidad"] * df["precio"]};
  % ---------- Panel B: pd.cut ----------
  \draw[panel] (8.8,0) rectangle (16.95,3.2);
  \node[tpanel] at (8.95,3.05) {B. De número a categoría};
  \fill[esumerTeal!45] (9.4,1.45) rectangle (11.7,2.0);
  \fill[esumerSky!45] (11.7,1.45) rectangle (14.0,2.0);
  \fill[esumerOrange!45] (14.0,1.45) rectangle (16.3,2.0);
  \node[font=\scriptsize\bfseries] at (10.55,1.725) {Joven};
  \node[font=\scriptsize\bfseries] at (12.85,1.725) {Adulto};
  \node[font=\scriptsize\bfseries] at (15.15,1.725) {Mayor};
  \node[font=\scriptsize\ttfamily] at (10.55,2.25) {(0, 25]};
  \node[font=\scriptsize\ttfamily] at (12.85,2.25) {(25, 45]};
  \node[font=\scriptsize\ttfamily] at (15.15,2.25) {(45, 120]};
  \foreach \x/\lim in {9.4/0, 11.7/25, 14.0/45, 16.3/120} {
    \draw[esumerNavy, line width=0.7pt] (\x,1.35) -- (\x,2.1);
    \node[font=\scriptsize] at (\x,1.1) {\lim};
  }
  \node[fn] at (12.9,0.45) {pd.cut(df["edad"], bins=[0, 25, 45, 120], labels=[...])};
  % ---------- Panel C: fechas ----------
  \draw[panel] (0.05,-3.6) rectangle (8.2,-0.4);
  \node[tpanel] at (0.15,-0.55) {C. Partes de una fecha};
  \node[dato, minimum width=2.0cm] (f) at (1.25,-2.1) {\textbf{fecha}\\2026-03-03};
  \node[font=\tiny, text=esumerNavyDark, align=center] at (1.25,-3.2) {FECHA\_REF =\\2026-09-30};
  \foreach \y/\fun/\res in {-1.3/{dt.to\_period("M")}/2026-03, -2.1/dt.quarter/1, -2.9/{(FECHA\_REF - fecha).dt.days}/211} {
    \draw[flecha] (f.east) -- (2.6,-2.1) |- (6.6,\y);
    \node[fn, anchor=south] at (4.6,\y) {\fun};
    \node[dato, minimum width=1.3cm, minimum height=0.55cm, fill=esumerTeal!25] at (7.35,\y) {\res};
  }
  % ---------- Panel D: map ----------
  \draw[panel] (8.8,-3.6) rectangle (16.95,-0.4);
  \node[tpanel] at (8.95,-0.55) {D. Traer un dato que conoces};
  \foreach \y/\pl/\val in {-1.35/Mensual/90.000, -2.0/Trimestral/240.000, -2.65/Anual/800.000} {
    \node[dato, minimum width=1.8cm, minimum height=0.5cm] (p) at (10.2,\y) {\pl};
    \node[dato, minimum width=1.6cm, minimum height=0.5cm, fill=esumerTeal!25] (v) at (12.9,\y) {\val};
    \draw[flecha] (p) -- (v);
  }
  \node[font=\scriptsize, text width=2.6cm, align=left] at (15.4,-2.0)
    {La lista de precios no venía en el CSV: la conoces por otra fuente y la traes con un diccionario.};
  \node[fn] at (12.9,-3.25) {df["plan"].map(\{"Mensual": 90000, "Anual": 800000\})};
\end{tikzpicture}

{\footnotesize\itshape Ejemplos de las Clases 8 y 9: la fecha y la
antigüedad son las de un socio real del dataset; los precios, los de la
lista del gimnasio.}
\end{center}

\begin{funciones}
Calcular una columna a partir de otras & \texttt{df["c"] = df["a"] * df["b"]} & 4, 8\\
Convertir un número en categorías & \texttt{pd.cut(s, bins=[...], labels=[...])} & 8\\
Convertir texto a fecha & \texttt{pd.to\_datetime(df["fecha"])} & 8, 9\\
Mes, trimestre, días transcurridos & \texttt{.dt.to\_period("M")}, \texttt{.dt.quarter}, \texttt{.dt.days} & 8, 9\\
Reemplazar cada valor según un diccionario & \texttt{s.map(\{...\})} & 8\\
Regla que usa varias columnas a la vez & \texttt{limpio.apply(funcion, axis=1)} & 8\\
\end{funciones}

\begin{codigo}[title=En tu notebook]
\begin{Verbatim}
FECHA_REF = pd.Timestamp("2026-09-30")          # fija: no uses la fecha de hoy

limpio["fecha"] = pd.to_datetime(limpio["fecha"])
limpio["mes"] = limpio["fecha"].dt.to_period("M")
limpio["antiguedad_dias"] = (FECHA_REF - limpio["fecha"]).dt.days

limpio["grupo_edad"] = pd.cut(limpio["edad"], bins=[0, 25, 45, 120],
                              labels=["Joven", "Adulto", "Mayor"])
limpio["grupo_edad"].value_counts()              # ¿quedó alguien sin grupo (NaN)?
\end{Verbatim}
\end{codigo}

\begin{cuidado}[Primero lo vectorizado]
\small Si existe una operación que trabaja sobre la columna completa
(\texttt{+}, \texttt{*}, \texttt{.str}, \texttt{.dt}, \texttt{pd.cut},
\texttt{.map}), úsala. Reserva \texttt{.apply(\ldots, axis=1)} para
reglas que necesitan varias columnas de la misma fila al mismo tiempo.
\end{cuidado}

\begin{checklist}
  \item Cada columna nueva responde a algo que necesito para mis preguntas.
  \item Revisé cada columna nueva con \texttt{value\_counts()} o
    \texttt{describe()}, y no generó nulos inesperados.
\end{checklist}

% ============================================================
\paso{5}{Integrar fuentes \normalsize(si tienes más de una)}{¿Cómo junto mis tablas sin perder ni duplicar filas en silencio?}

Si todos tus datos están en una sola tabla, puedes saltar este paso. Si
tienes un catálogo (productos, sedes, planes) o archivos que llegan por
partes, la integración tiene cuatro momentos, y el tercero es el único que
escribe la unión.

\begin{center}
\begin{tikzpicture}[
  momento/.style={caja, minimum width=3.7cm, text width=3.4cm, minimum height=2.1cm, fill=esumerBg},
]
  \foreach \i/\x/\tit/\fun in {
      1/2.2/{Revisa la llave}/{is\_unique \textperiodcentered\ dtype\\set(a) - set(b)},
      2/6.5/{Normaliza la llave}/{str.strip()\\replace()},
      3/10.8/{Une}/{merge(on=..., how=...,\\validate=...)},
      4/15.1/{Verifica}/{indicator=True\\len antes y después}} {
    \node[momento] (m\i) at (\x,0) {{\small\bfseries \tit}\\[4pt]{\scriptsize\ttfamily \fun}};
    \node[num] at (m\i.north west) {\i};
  }
  \foreach \i [evaluate=\i as \j using int(\i+1)] in {1,2,3}
    {\draw[flecha] (m\i) -- (m\j);}
\end{tikzpicture}
\end{center}

\textbf{El parámetro \texttt{how} decide qué filas sobreviven} (verde: se
conserva; gris: se descarta). En la mayoría de los análisis quieres
\texttt{how="left"}: partir de tu tabla principal y agregarle información
sin perder a nadie.

\begin{center}
\begin{tikzpicture}
  \foreach \k/\h/\desc in {0/inner/solo lo que empareja, 1/left/toda la tabla izquierda, 2/right/toda la tabla derecha, 3/outer/{todo, de ambas}} {
    \begin{scope}[xshift=\k*4.3cm]
      \coordinate (A) at (1.65,0);
      \coordinate (B) at (2.55,0);
      \fill[esumerGray] (A) circle (0.8);
      \fill[esumerGray] (B) circle (0.8);
      \ifnum\k=1 \fill[esumerTeal] (A) circle (0.8); \fi
      \ifnum\k=2 \fill[esumerTeal] (B) circle (0.8); \fi
      \ifnum\k=3 \fill[esumerTeal] (A) circle (0.8); \fill[esumerTeal] (B) circle (0.8); \fi
      \begin{scope}
        \clip (A) circle (0.8);
        \fill[esumerTeal] (B) circle (0.8);
      \end{scope}
      \draw[esumerNavy, line width=0.8pt] (A) circle (0.8);
      \draw[esumerNavy, line width=0.8pt] (B) circle (0.8);
      \node[font=\scriptsize] at (1.2,1.05) {izquierda};
      \node[font=\scriptsize] at (3.0,1.05) {derecha};
      \node[font=\small\ttfamily] at (2.1,-1.15) {how="\h"};
      \node[font=\scriptsize\itshape] at (2.1,-1.6) {\desc};
    \end{scope}
  }
\end{tikzpicture}
\end{center}

\begin{funciones}
Comprobar que la llave no se repite en el catálogo & \texttt{tabla["ciudad"].is\_unique} & 10\\
Ver valores sin pareja antes de unir & \texttt{set(limpio["ciudad"]) - set(tabla["ciudad"])} & 10\\
Unir por una llave, validando la relación & \texttt{limpio.merge(tabla, on="ciudad", how="left", validate="many\_to\_one")} & 10\\
Detectar lo que quedó sin pareja & \texttt{merge(..., indicator=True)} y \texttt{["\_merge"].value\_counts()} & 10\\
Apilar tablas con las mismas columnas & \texttt{pd.concat([a, b], ignore\_index=True)} & 10\\
\end{funciones}

\begin{codigo}[title=En tu notebook]
\begin{Verbatim}
print(tabla["ciudad"].is_unique)
print(set(limpio["ciudad_limpia"]) - set(tabla["ciudad"]))    # debería quedar vacío

unida = limpio.merge(tabla, left_on="ciudad_limpia", right_on="ciudad",
                     how="left", validate="many_to_one", indicator=True)
print(unida["_merge"].value_counts())        # ¿cuántos left_only?
print(len(limpio), "->", len(unida))        # en un left many_to_one deben ser iguales
\end{Verbatim}
\end{codigo}

\begin{recuerda}
\small En la Clase 10, la sede de Itagüí estaba escrita \texttt{"Itagui"}
y sus socios quedaban sin sede sin ningún error; al apilar las nuevas
inscripciones, el \texttt{id\_socio} 7 apareció dos veces con dos personas
distintas; y una sede repetida en el catálogo convirtió 50 socios en 72
filas. Las tres cosas se detectan con los chequeos de esta página.
\end{recuerda}

\begin{checklist}
  \item La tabla unida tiene el número de filas que esperaba, y sé por qué.
  \item Revisé los \texttt{left\_only} y \texttt{right\_only}, y anoté qué hice con ellos.
  \item Usé \texttt{validate} en cada \texttt{merge}.
\end{checklist}

% ============================================================
\paso{6}{Agrupar y resumir}{¿Cómo se comporta cada grupo, comparado con los demás?}

\texttt{groupby} hace siempre lo mismo, en tres tiempos: divide la tabla en
grupos, aplica un cálculo a cada uno y junta los resultados en una tabla
nueva, con una fila por grupo.

\begin{center}
\begin{tikzpicture}[
  c/.style={draw=esumerGray!70!black, line width=0.4pt, minimum height=0.5cm, inner sep=1pt, font=\scriptsize},
  h/.style={c, fill=esumerNavy, text=white, font=\scriptsize\bfseries},
  et/.style={font=\small\bfseries, text=esumerNavy},
]
  % colores por grupo
  \colorlet{gM}{esumerTeal!30}
  \colorlet{gB}{esumerSky!30}
  \colorlet{gE}{esumerOrange!35}
  % ---------- Tabla original ----------
  \node[et] at (1.7,0.95) {Tabla};
  \node[h, minimum width=2.0cm] at (1.0,0.25) {ciudad};
  \node[h, minimum width=1.4cm] at (2.7,0.25) {edad};
  \foreach \i/\ci/\e/\g in {1/Medellín/30/gM, 2/Bogotá/50/gB, 3/Envigado/20/gE,
                            4/Medellín/40/gM, 5/Bogotá/30/gB, 6/Envigado/40/gE} {
    \node[c, minimum width=2.0cm, fill=\g] at (1.0,0.25-0.5*\i) {\ci};
    \node[c, minimum width=1.4cm, fill=\g] at (2.7,0.25-0.5*\i) {\e};
  }
  % ---------- Dividir ----------
  \node[et] at (6.6,0.95) {1. Dividir};
  \node[fn] at (6.6,0.45) {groupby("ciudad")};
  \foreach \ci/\g/\ya/\e/\f in {Medellín/gM/-0.25/30/40, Bogotá/gB/-1.5/50/30, Envigado/gE/-2.75/20/40} {
    \node[c, minimum width=2.0cm, fill=\g] at (6.0,\ya) {\ci};
    \node[c, minimum width=1.2cm, fill=\g] at (7.6,\ya) {\e};
    \node[c, minimum width=2.0cm, fill=\g] at (6.0,\ya-0.5) {\ci};
    \node[c, minimum width=1.2cm, fill=\g] at (7.6,\ya-0.5) {\f};
  }
  \draw[flecha] (3.6,-1.5) -- (4.8,-1.5);
  % ---------- Aplicar ----------
  \node[et] at (10.5,0.95) {2. Aplicar};
  \node[fn] at (10.5,0.45) {["edad"].mean()};
  \foreach \ya/\val/\g in {-0.5/35/gM, -1.75/40/gB, -3.0/30/gE} {
    \draw[flecha] (8.35,\ya) -- (9.9,\ya);
    \node[circle, draw=esumerNavy, fill=\g, minimum size=0.75cm, font=\small\bfseries] at (10.5,\ya) {\val};
  }
  % ---------- Combinar ----------
  \node[et] at (14.85,0.95) {3. Combinar};
  \node[fn] at (14.85,0.45) {una fila por grupo};
  \node[h, minimum width=2.0cm] at (14.0,-1.0) {ciudad};
  \node[h, minimum width=1.7cm] at (15.85,-1.0) {edad prom.};
  \foreach \i/\ci/\val/\g in {1/Medellín/35/gM, 2/Bogotá/40/gB, 3/Envigado/30/gE} {
    \node[c, minimum width=2.0cm, fill=\g] at (14.0,-1.0-0.5*\i) {\ci};
    \node[c, minimum width=1.7cm, fill=\g] at (15.85,-1.0-0.5*\i) {\val};
  }
  \foreach \ya/\yb in {-0.5/-1.5, -1.75/-2.0, -3.0/-2.5}
    {\draw[flecha] (10.9,\ya) -- (12.95,\yb);}
\end{tikzpicture}

{\footnotesize\itshape Datos ilustrativos (seis filas) para mostrar la mecánica.}
\end{center}

\begin{funciones}
Contar filas por grupo & \texttt{limpio.groupby("ciudad").size()} & 9\\
Varios indicadores a la vez, con nombre & \texttt{.agg(clientes=("id\_cliente", "count"), monto=("monto", "sum"))} & 9\\
Ordenar el resumen & \texttt{.sort\_values("clientes", ascending=False)} & 9\\
Tabla cruzada de dos categorías & \texttt{pd.crosstab(limpio["ciudad"], limpio["plan"])} & 9\\
Total del grupo en cada fila (para participaciones) & \texttt{.groupby("ciudad")["monto"].transform("sum")} & 9\\
Resumen por mes & \texttt{limpio.groupby("mes").size()} & 9\\
\end{funciones}

\begin{codigo}[title=En tu notebook]
\begin{Verbatim}
resumen = (limpio.groupby("ciudad_limpia")
                 .agg(clientes=("id_cliente", "count"),
                      monto_total=("monto", "sum"),
                      edad_promedio=("edad", "mean"))
                 .sort_values("clientes", ascending=False)
                 .round(1))
resumen["participacion"] = (resumen["clientes"] / resumen["clientes"].sum() * 100).round(1)
resumen
\end{Verbatim}
\end{codigo}

\begin{cuidado}[Un promedio sin su conteo engaña]
\small Incluye siempre una columna con el número de filas de cada grupo.
En la Clase 9, Itagüí tenía solo 2 socios: su promedio se veía igual de
``serio'' que el de Medellín, con 22. Si un grupo tiene muy pocas filas,
dilo junto al resultado o júntalo con otros.
\end{cuidado}

\begin{checklist}
  \item El resumen responde directamente a una de mis preguntas del Paso 0.
  \item Tiene una columna de conteo, está ordenado y redondeado.
  \item La suma de los conteos coincide con el total de filas de \texttt{limpio}.
\end{checklist}

% ============================================================
\paso{7}{Visualizar y comunicar}{¿Qué gráfico responde mi pregunta, y qué conclusión quiero que se lleve quien lo vea?}

Vuelve a tus preguntas del Paso 0: el tipo de pregunta decide el gráfico.
Todos se hacen con el método \texttt{.plot} de pandas, directamente sobre
la columna o la tabla resumen.

\begin{center}
\begin{tikzpicture}[
  rama/.style={draw=none, rounded corners=3pt, text=white, font=\small\bfseries,
               minimum width=3.8cm, text width=3.6cm, align=center, minimum height=1.25cm, inner sep=3pt},
  opc/.style={caja, minimum width=3.8cm, text width=3.6cm, align=left, font=\scriptsize, inner sep=4pt},
]
  \node[caja, fill=esumerNavy, text=white, font=\small\bfseries, text width=7cm] (raiz) at (8.65,0)
    {¿Qué pregunta responde tu gráfico?};
  \node[rama, fill=esumerTeal] (r1) at (2.2,-1.7) {Comparación\\[-1pt]{\scriptsize\mdseries ¿cuál es mayor?}};
  \node[rama, fill=esumerSky] (r2) at (6.5,-1.7) {Composición\\[-1pt]{\scriptsize\mdseries ¿cómo se reparte?}};
  \node[rama, fill=esumerOrange] (r3) at (10.8,-1.7) {Distribución\\[-1pt]{\scriptsize\mdseries ¿cómo se dispersa?}};
  \node[rama, fill=esumerMagenta] (r4) at (15.1,-1.7) {Relación\\[-1pt]{\scriptsize\mdseries ¿se mueven juntas?}};
  \node[opc, anchor=north] at (2.2,-2.55)
    {\textbf{Pocas categorías:} barras\\{\ttfamily s.plot.barh()}\\[3pt]
     \textbf{En el tiempo:} línea\\{\ttfamily s.plot(marker="o")}};
  \node[opc, anchor=north] at (6.5,-2.55)
    {\textbf{Entre grupos:} apiladas\\{\ttfamily t.plot.bar(stacked=True)}\\[3pt]
     \textbf{3 a 5 partes:} pastel\\{\ttfamily s.plot.pie()}};
  \node[opc, anchor=north] at (10.8,-2.55)
    {\textbf{Una variable:} histograma\\{\ttfamily s.plot.hist()}\\[3pt]
     \textbf{Por grupos:} boxplot\\{\ttfamily df.boxplot(column="edad",\\\hspace*{1.5em}by="grupo")}};
  \node[opc, anchor=north] at (15.1,-2.55)
    {\textbf{Dos numéricas:} dispersión\\{\ttfamily df.plot.scatter(x="edad",\\\hspace*{1.5em}y="monto")}\\[3pt]
     \textbf{Acompáñala con:}\\{\ttfamily df["a"].corr(df["b"])}};
  \foreach \i in {1,2,3,4} {\draw[flecha] (raiz.south) -- ++(0,-0.35) -| (r\i.north);}
\end{tikzpicture}

{\footnotesize\itshape \texttt{s}: una columna o un resumen (\texttt{Series});
\texttt{t}: una tabla cruzada; \texttt{df}: tu tabla. Ejemplos completos en el notebook de la Clase 11.}
\end{center}

\begin{codigo}[title=En tu notebook: un gráfico completo]
\begin{Verbatim}
import matplotlib.pyplot as plt

resumen["clientes"].sort_values().plot.barh(color="#0CB3B6",
                                            title="Escribe aquí tu conclusión, no una descripción")
plt.xlabel("Clientes")
plt.show()
\end{Verbatim}
\end{codigo}

\textbf{Antes de mostrarlo} (la lista de la Clase 11):

\begin{checklist}
  \item El título dice la conclusión (``Medellín concentra el 44\% de los
    socios''), no la descripción (``Socios por ciudad'').
  \item En barras, el eje empieza en cero.
  \item Los ejes tienen nombre y unidad, y las categorías están ordenadas por valor.
  \item Cada color significa algo; no hay rejillas, leyendas ni bordes que sobren.
  \item Sin doble eje: si son dos variables con unidades distintas, van en dos gráficos.
\end{checklist}

\textbf{Del gráfico al hallazgo.} Un gráfico no es todavía un hallazgo. El
hallazgo es la frase que alguien que toma decisiones puede usar:

\begin{center}
\begin{tikzpicture}[
  pieza/.style={caja, minimum width=3.8cm, text width=3.5cm, minimum height=1.9cm},
  tit/.style={font=\bfseries\small, text=white, fill=esumerTeal, rounded corners=2pt,
              minimum width=3.8cm, inner sep=3pt},
]
  \foreach \x/\tit in {2.2/Evidencia, 6.5/Gráfico, 10.8/Hallazgo, 15.1/Recomendación}
    {\node[tit] at (\x,1.3) {\tit};}
  \node[pieza] (a) at (2.2,0) {\small 22 de los 50 socios están en Medellín};
  \node[pieza] (b) at (6.5,0) {\small Barras horizontales, ordenadas, eje desde cero};
  \node[pieza] (c) at (10.8,0) {\small Medellín concentra el 44\% de los socios};
  \node[pieza] (d) at (15.1,0) {\small Priorizar allí la campaña de renovación};
  \draw[flecha] (a) -- (b);
  \draw[flecha] (b) -- (c);
  \draw[flecha] (c) -- (d);
\end{tikzpicture}

{\footnotesize\itshape Cifras reales del dataset limpio del gimnasio; la recomendación es un ejemplo.}
\end{center}

% ============================================================
\clearpage
\section*{\color{esumerNavy}Entregable de la sesión}

Al final de la clase, comparte tu notebook y esta página completa. Es
también el punto de partida de la Clase 13: el archivo limpio que guardes
hoy es el que vas a cargar en Power BI.

\begin{codigo}[title=Última celda de tu notebook]
\begin{Verbatim}
limpio.to_csv("mis_datos_limpios.csv", index=False)    # descárgalo desde el panel de archivos
\end{Verbatim}
\end{codigo}

\textbf{Mis hallazgos}

\renewcommand{\arraystretch}{3.2}
\begin{tabularx}{\textwidth}{@{}>{\centering\arraybackslash}p{0.5cm} >{\raggedright\arraybackslash}X >{\raggedright\arraybackslash}X >{\raggedright\arraybackslash}X@{}}
\toprule
& \textbf{Evidencia (cifra y gráfico)} & \textbf{Hallazgo} & \textbf{Recomendación}\\
\midrule
1 & & & \\ \hline
2 & & & \\ \hline
3 & & & \\
\bottomrule
\end{tabularx}\par
\renewcommand{\arraystretch}{1.15}

\vspace{8pt}
\textbf{Bitácora de decisiones}

Una fila por cada problema que encontraste. La primera es un ejemplo del
curso; quien lea tu análisis debe poder entender, solo con esta tabla, qué
cambiaste en los datos y por qué.

\renewcommand{\arraystretch}{2.0}
\begin{tabularx}{\textwidth}{@{}>{\centering\arraybackslash}p{0.9cm} >{\raggedright\arraybackslash}X >{\raggedright\arraybackslash}X >{\raggedright\arraybackslash}p{2.0cm} >{\raggedright\arraybackslash}X@{}}
\toprule
\textbf{Paso} & \textbf{Qué encontré} & \textbf{Qué decidí} & \textbf{Filas afectadas} & \textbf{Por qué}\\
\midrule
{\small\itshape 3} & {\small\itshape \texttt{edad} tenía texto (\texttt{"treinta"}) y celdas vacías} & {\small\itshape Convertir con \texttt{to\_numeric} e imputar con la mediana} & {\small\itshape las que no eran un número} & {\small\itshape La mediana no se deja arrastrar por valores extremos}\\ \hline
& & & & \\ \hline
& & & & \\ \hline
& & & & \\ \hline
& & & & \\ \hline
& & & & \\
\bottomrule
\end{tabularx}\par
\renewcommand{\arraystretch}{1.15}

\end{document}
